\chapter{Classical Machine Learning Foundations}
\label{ch:classical}

The comparison at the heart of this thesis is only as strong as its classical
side. This chapter fixes the supervised-learning formalism and notation, then
reviews the model families that serve as baselines --- with particular depth on
kernel methods, since the mathematics of Chapter~\ref{ch:qml}'s quantum kernels
is the same mathematics --- and closes with the evaluation methodology that the
fairness protocol of Chapter~\ref{ch:methodology} operationalizes. Standard
references for this material are \citet{bishop2006prml} and
\citet{hastie2009esl}.

\section{Supervised learning}
\label{sec:ml-setup}

In supervised binary classification we observe a sample
$\data = \{(\bm{x}_i, y_i)\}_{i=1}^{m}$ drawn i.i.d.\ from an unknown
distribution $\Dcal$ over $\mathcal{X} \times \mathcal{Y}$ with
$\mathcal{X} \subseteq \R^d$ and $\mathcal{Y} = \{0, 1\}$. A learning
algorithm selects a hypothesis $h \colon \mathcal{X} \to \mathcal{Y}$ (often
via a real-valued score $f$ with $h = \mathbf{1}[f \ge \tau]$) from a
hypothesis class $\Fcal$, seeking small \emph{risk}
\begin{equation}
  R(h) = \E_{(\bm{x}, y) \sim \Dcal} \bigl[ \ell(h(\bm{x}), y) \bigr],
\end{equation}
for a loss $\ell$, while only the \emph{empirical risk}
$\widehat{R}(h) = \tfrac{1}{m} \sum_i \ell(h(\bm{x}_i), y_i)$ is observable.
The gap between the two is the object of generalization theory; its practical
management --- model selection without contaminating the estimate of $R$ ---
is the subject of Section~\ref{sec:ml-evaluation}.

The classical decomposition of expected error into \emph{bias},
\emph{variance}, and irreducible noise \citep{hastie2009esl} shapes intuition
throughout this thesis: rich hypothesis classes (deep trees, wide networks,
expressive quantum circuits) reduce bias at the price of variance, and
regularization --- explicit penalties, restricted architectures, or early
stopping --- trades back along the same curve. Two caveats are important.
Modern over-parameterized models can defy the naive U-shaped picture
\citep{goodfellow2016deep}, and the same caution applies to quantum models:
expressivity arguments alone predict neither good nor bad generalization
\citep{gilfuster2024rethinking}, which is why this thesis measures rather than
assumes (RQ2, RQ5).

\section{Linear models and logistic regression}
\label{sec:ml-linear}

The simplest score is affine, $f(\bm{x}) = \bm{w}^{\top} \bm{x} + b$.
Logistic regression passes it through the sigmoid
$\sigma(t) = (1 + e^{-t})^{-1}$, interprets
$\sigma(f(\bm{x}))$ as $\Pr(y = 1 \mid \bm{x})$, and fits $(\bm{w}, b)$ by
minimizing the $\ell_2$-regularized negative log-likelihood
\begin{equation}
  \min_{\bm{w}, b} \; \frac{1}{C} \cdot \frac{\norm{\bm{w}}^2}{2}
  + \sum_{i=1}^{m} \log\!\left( 1 + e^{-\tilde{y}_i (\bm{w}^{\top}\bm{x}_i + b)} \right),
  \qquad \tilde{y}_i = 2y_i - 1 \in \{\pm 1\},
\end{equation}
a smooth convex problem. Logistic regression contributes two things to this
thesis: a floor that any nonlinear model must clear to justify its complexity,
and a diagnostic --- tasks on which it performs at chance (the concentric
circles and sign-parity constructions of Chapter~\ref{ch:methodology}) are
certified to require nonlinear structure.

\section{Kernel methods and support vector machines}
\label{sec:ml-kernels}

Kernel methods make nonlinearity affordable by moving it into an inner
product. A \emph{feature map} $\phi \colon \mathcal{X} \to \Hil$ into a
(possibly infinite-dimensional) Hilbert space defines the \emph{kernel}
\begin{equation}
  \kappa(\bm{x}, \bm{x}') = \braket{\phi(\bm{x})}{\phi(\bm{x}')}_{\Hil},
  \label{eq:kernel}
\end{equation}
which is symmetric and positive semidefinite; conversely (Mercer), any such
kernel is an inner product in \emph{some} feature space. Algorithms that touch
data only through inner products can substitute $\kappa$ and operate in
$\Hil$ implicitly --- the \emph{kernel trick}.

The support vector machine (SVM) \citep{cortes1995svm} seeks the
maximum-margin separator in $\Hil$. Its soft-margin primal,
\begin{equation}
  \min_{\bm{w}, b, \bm{\xi}} \;\; \frac{1}{2} \norm{\bm{w}}^2
    + C \sum_{i=1}^{m} \xi_i
  \quad \text{s.t.} \quad
  \tilde{y}_i \bigl( \braket{\bm{w}}{\phi(\bm{x}_i)} + b \bigr) \ge 1 - \xi_i,
  \;\; \xi_i \ge 0,
\end{equation}
has the dual
\begin{equation}
  \max_{\bm{\alpha}} \; \sum_{i=1}^{m} \alpha_i
    - \frac{1}{2} \sum_{i,j=1}^{m} \alpha_i \alpha_j \tilde{y}_i \tilde{y}_j\,
      \kappa(\bm{x}_i, \bm{x}_j)
  \quad \text{s.t.} \quad
  0 \le \alpha_i \le C, \;\; \textstyle\sum_i \alpha_i \tilde{y}_i = 0,
  \label{eq:svm-dual}
\end{equation}
with decision function
$f(\bm{x}) = \sum_i \alpha_i \tilde{y}_i \kappa(\bm{x}_i, \bm{x}) + b$
supported on the training points with $\alpha_i > 0$ (the \emph{support
vectors}). Two properties of \eqref{eq:svm-dual} matter here. First, the data
enter only through the \emph{Gram matrix}
$K_{ij} = \kappa(\bm{x}_i, \bm{x}_j)$ --- so any procedure that can estimate
pairwise kernel values, including a quantum computer estimating state
fidelities, can drive an otherwise entirely classical SVM. This is exactly the
quantum-kernel pipeline of Section~\ref{sec:qml-kernels}, and it is why the
SVM receives the most detailed treatment of any classical model in this
thesis. Second, the optimization is convex: given the Gram matrix, training
has no local-minimum or initialization pathologies, in sharp contrast to the
variational models of Section~\ref{sec:qml-vqa}.

The canonical general-purpose kernel is the radial basis function (RBF),
\begin{equation}
  \kappa_{\mathrm{RBF}}(\bm{x}, \bm{x}')
    = \exp\bigl( -\gamma \norm{\bm{x} - \bm{x}'}^2 \bigr),
  \label{eq:rbf}
\end{equation}
whose bandwidth $\gamma$ interpolates between nearly linear behavior
($\gamma \to 0$) and nearest-neighbor-like memorization ($\gamma \to \infty$).
The RBF-SVM is the single most important classical baseline in this thesis:
it is the like-for-like classical counterpart of the quantum-kernel SVM ---
same learning algorithm, different similarity function --- so any performance
difference is attributable to the kernel itself.

\section{Neural networks and gradient-based training}
\label{sec:ml-nn}

A multilayer perceptron (MLP) with one hidden layer of width $k$ computes
\begin{equation}
  f(\bm{x}) = \bm{w}_2^{\top}\, g\bigl( W_1 \bm{x} + \bm{b}_1 \bigr) + b_2,
\end{equation}
with an elementwise nonlinearity $g$ (ReLU in this thesis), giving
$dk + 2k + 1$ trainable parameters. Universal-approximation results guarantee
expressivity in the wide limit; what is actually achievable is decided by
optimization --- non-convex, performed by stochastic gradient descent on
mini-batches with gradients from backpropagation, in this work using the Adam
optimizer \citep{kingma2015adam} --- and by regularization ($\ell_2$ weight
penalties here) \citep{goodfellow2016deep}.

The MLP has a designated role in the fairness protocol: it is the
\emph{parameter-matched} classical mirror of the variational quantum
classifier. A VQC with $p$ trainable rotation angles is compared against MLPs
constrained to approximately $p$ parameters, so that RQ3's hybrid-vs-classical
comparisons oppose architectures of comparable capacity rather than comparable
marketing. Both model families also share an optimization story --- noisy
gradient estimates of a non-convex objective --- which makes their trainability
comparison (barren plateaus versus vanishing gradients,
Section~\ref{sec:qml-trainability}) meaningful.

\section{Ensembles: random forests and gradient boosting}
\label{sec:ml-ensembles}

Decision-tree ensembles are the strongest general-purpose baselines on small-%
to-medium tabular data --- precisely this thesis's data regime --- and their
absence is a recurring flaw in QML comparisons. A \emph{random forest}
\citep{breiman2001rf} averages trees grown on bootstrap resamples with random
feature subsetting at each split, reducing variance while keeping bias low.
\emph{Gradient boosting} \citep{chen2016xgboost} instead builds trees
sequentially, each fit to the gradient of the loss at the current ensemble's
predictions, with shrinkage and subsampling controlling capacity; this work
uses the XGBoost implementation. Neither model consumes the $[0,\pi]$ feature
scaling in any essential way --- trees are invariant to monotone feature maps
--- which makes them a useful control: a model family whose performance is
unaffected by the encoding-motivated preprocessing shared across the
comparison.

\section{Instance-based learning}
\label{sec:ml-knn}

$k$-nearest-neighbors classifies by majority vote among the $k$ closest
training points under a Minkowski metric. It is included as the minimal
``pure memorization'' baseline: no training, no parametric structure, purely
local generalization. Its behavior on the engineered synthetic tasks is
diagnostic --- it excels when local neighborhoods are label-pure (moons,
circles) and degrades when they are not (the sign-parity construction, where
distractor dimensions corrupt the metric) --- providing a reference point that
separates ``the task needs nonlinearity'' from ``the task needs
\emph{non-local} structure.''

\section{Evaluation methodology}
\label{sec:ml-evaluation}

\paragraph{Metrics.}
Accuracy is the headline metric on balanced tasks but degrades to
uninformativeness under imbalance. This thesis therefore reports four
complementary quantities per evaluation: \emph{accuracy};
\emph{balanced accuracy}, the mean of per-class recalls, robust to the
moderate imbalance of several included datasets; $F_1$, the harmonic mean of
precision and recall for the positive class; and \emph{ROC-AUC}, the
probability that a random positive instance outscores a random negative one,
which evaluates the score function independently of a threshold.

\paragraph{Resampling.}
A single train/test split yields a point estimate whose variance is unknown.
The protocol here uses $k$-fold \emph{stratified} cross-validation ($k = 5$),
which preserves class proportions per fold, repeated over five different
shuffling seeds --- 25 evaluations per model--dataset pair --- reporting means
and standard deviations across folds. Crucially, the split indices are
generated \emph{once}, stored, and reused by every model ever compared:
resampling variance then cancels in model comparisons, turning the folds into
matched pairs.

\paragraph{Hyperparameter optimization without leakage.}
Hyperparameters must be selected using training data only; selecting them on
test folds is the most common form of optimistic bias. The protocol fixes a
matched budget --- one 50-trial Optuna study \citep{akiba2019optuna} per
model--dataset pair, using the tree-structured Parzen estimator sampler with a
fixed seed, each trial scored by an internal 3-fold cross-validation confined
to the training portion of one designated outer fold --- after which the
selected configuration is frozen and evaluated on all 25 outer folds. The
residual weakness of this design (the tuning fold's training data reappears in
other folds' evaluations) is deliberate, tractability-motivated, and assessed
in Chapter~\ref{ch:discussion}; its key property is \emph{symmetry}: every
model, classical or quantum, is tuned by exactly the same procedure with
exactly the same budget \citep{bowles2024better}.

\paragraph{Statistical comparison.}
With 25 paired evaluations per model pair, per-dataset differences are tested
with the paired Wilcoxon signed-rank test, and comparisons across the whole
dataset suite follow the rank-based methodology of \citet{demsar2006comparisons},
including critical-difference diagrams for the headline comparison.
Significance levels are stated in advance in Chapter~\ref{ch:methodology};
non-significant differences are reported as such rather than narrativized.
